\documentclass[aspectratio=169]{beamer}

% --- PAQUETS ---
\usepackage[utf8]{inputenc}
\usepackage[T1]{fontenc}
\usepackage[french]{babel}
\usepackage{graphicx}
\usepackage{listings}
\usepackage{xcolor}
\usepackage{tikz}
\usetikzlibrary{shapes,arrows,positioning,fit,backgrounds,calc}

% --- THÈME ---
\usetheme{Madrid}
\usecolortheme{whale}

% --- CONFIGURATION CODE ---
\definecolor{codegreen}{rgb}{0,0.6,0}
\definecolor{codegray}{rgb}{0.5,0.5,0.5}
\definecolor{codepurple}{rgb}{0.58,0,0.82}
\definecolor{backcolour}{rgb}{0.95,0.95,0.92}

\lstdefinestyle{mystyle}{
    backgroundcolor=\color{backcolour},   
    commentstyle=\color{codegreen},
    keywordstyle=\color{magenta},
    numberstyle=\tiny\color{codegray},
    stringstyle=\color{codepurple},
    basicstyle=\ttfamily\scriptsize,
    breakatwhitespace=false,         
    breaklines=true,                 
    captionpos=b,                    
    keepspaces=true,                 
    numbers=left,                    
    numbersep=5pt,                  
    showspaces=false,                
    showstringspaces=false,
    showtabs=false,                  
    tabsize=2
}
\lstset{style=mystyle}

% --- INFOS PRÉSENTATION ---
\title[Big Data]{Mise en place d'une architecture distribuée\\pour l'analyse de logs web}
\subtitle{HDFS, Batch, Streaming et NoSQL avec Docker}
\author{Matthieu PELINGRE} 
\institute{IDMC / M2SDD}
\date{7 février 2026}

\begin{document}

% --- SLIDE 1: TITRE ---
\begin{frame}
  \titlepage
\end{frame}

% --- SLIDE 2: PLAN ---
\begin{frame}{Plan de la présentation}
  \tableofcontents
\end{frame}

% --- SECTION 1: CONTEXTE ---
\section{Contexte et objectifs}

\begin{frame}{Contexte et problématique}
  \begin{columns}
    \column{0.6\textwidth}
    \textbf{Scénario :}
    \begin{itemize}
        \item Site e-commerce (Cosmétiques).
        \item Fort volume de trafic générant des logs serveurs (\texttt{web\_server.log}).
    \end{itemize}
    
    \vspace{0.5cm}
    \textbf{Besoin :}
    \begin{itemize}
        \item Transformer des données brutes en informations décisionnelles.
        \item Besoin de scalabilité (Big Data) et de réactivité (Temps réel).
    \end{itemize}

    \column{0.4\textwidth}
    \begin{alertblock}{Objectifs Techniques}
        \begin{enumerate}
            \item \textbf{Stocker} (HDFS).
            \item \textbf{Traiter} (Spark Batch).
            \item \textbf{Surveiller} (Spark Streaming).
            \item \textbf{Persister} (MongoDB).
        \end{enumerate}
    \end{alertblock}
  \end{columns}
\end{frame}

% --- SECTION 2: ARCHITECTURE ---
\section{Architecture technique}

\begin{frame}{Architecture déployée (Docker)}
    \begin{center}
    \resizebox{\textwidth}{!}{
    \begin{tikzpicture}[
        node distance=1.5cm and 2cm,
        box/.style={rectangle, draw, rounded corners, fill=blue!10, align=center, minimum height=1.2cm, minimum width=2.5cm, font=\footnotesize},
        database/.style={cylinder, shape border rotate=90, draw, fill=green!10, aspect=0.25, align=center, minimum height=1.5cm, minimum width=1.5cm, font=\footnotesize},
        source/.style={rectangle, draw, dashed, fill=yellow!10, align=center, minimum height=1cm, minimum width=2.5cm, font=\footnotesize},
        arrow/.style={->, >=stealth, thick, color=black!70}
    ]

    % --- Row 1: Stream Flow ---
    \node[source] (gen) {Générateur Python\\(socket stream)};
    \node[box, right=of gen, fill=blue!20] (spark) {\textbf{Apache Spark}\\(Master + 2 Workers)};
    \node[database, right=of spark, fill=orange!20] (mongo) {\textbf{MongoDB}\\(Logs DB)};

    % --- Row 2: Batch Flow ---
    \node[source, below=1.5cm of gen] (file) {Fichier Logs\\(données brutes)};
    \node[database, right=of file] (hdfs) {\textbf{Hadoop HDFS}\\(NN + 2 DN)};

    % --- Connections ---
    % Stream path
    \draw[arrow] (gen) -- node[above, font=\tiny, color=red] {TCP :9998} (spark);
    \draw[arrow] (spark) -- node[above, font=\tiny] {Documents} (mongo);

    % Batch path
    \draw[arrow] (file) -- node[above, font=\tiny] {Upload (Client)} (hdfs);
    \draw[arrow] (hdfs) -- node[right, font=\tiny, pos=0.4] {Lecture Batch} (spark);

    % --- Grouping ---
    \begin{pgfonlayer}{background}
        \node[fit=(spark) (hdfs) (mongo), draw, dashed, red!50, fill=red!5, inner sep=0.5cm, rounded corners, label={[anchor=south east, red, font=\scriptsize]south east:Network}] {};
    \end{pgfonlayer}

    \end{tikzpicture}
    }
    \end{center}
\end{frame}

\begin{frame}{Stratégie d'ingestion \& développement}
  \begin{block}{Automatisation}
    \begin{itemize}
        \item \textbf{Batch :} conteneur \texttt{hdfs-client} éphémère. Attend la fin du démarrage du NameNode, puis téléverse les données.
        \item \textbf{Stream :} service \texttt{data-generator} résilient (boucle infinie + reconnexion automatique).
    \end{itemize}
  \end{block}

  \begin{block}{Qualité de code}
    \begin{itemize}
        \item Création d'un module \texttt{utils.py} partagé.
        \item Centralise la logique de parsing : Regex pour CLF (\textit{Common Log Format}).
        \item Distribué aux workers via l'option \texttt{--py-files}.
    \end{itemize}
  \end{block}
\end{frame}

% --- SECTION 3: BATCH ---
\section{Traitements Batch (le passé)}

\begin{frame}[fragile]{Analyse Batch : méthodologie}
    \textbf{Approche MapReduce avec PySpark}
    \begin{itemize}
        \item Lecture directe depuis HDFS : \texttt{hdfs://namenode:9000/...}
        \item Parsing des lignes via \texttt{utils.parse\_log\_line()} $\rightarrow$ RDD
    \end{itemize}
    
    \vspace{0.1cm}
    
    \begin{columns}[t]
        \column{0.5\textwidth}
        \textbf{A. Répartition codes HTTP}
        \begin{itemize}
            \item RDD en cache pour :
            \begin{itemize}
                \item Comptage total.
                \item Comptage par code HTTP.
            \end{itemize}
            \item Conversion du RDD en DataFrame.
            \item \textbf{Enrichissement :} Jointure avec descriptions.
            \item Calcul de fréquence.
            \item Sauvegarde MongoDB (overwrite).
        \end{itemize}

        \column{0.5\textwidth}
        \textbf{B. Top 10 IPs}
        \begin{itemize}
            \item Comptage par adresse IP.
            \item Utilisation de \texttt{takeOrdered(10)}.
            \item Conversion du RDD en DataFrame.
            \item Sauvegarde MongoDB (overwrite).
        \end{itemize}
    \end{columns}
\end{frame}

\begin{frame}[fragile]{Résultats Batch (codes HTTP)}
    
    \begin{exampleblock}{Affichage console (Master Spark)}
\begin{lstlisting}[language=Bash]
Demarrage du traitement Batch : AnalyzeLogsCodesHTTP
Total des requetes analysees : 4369067
Resultats : Frequence des codes HTTP
+----+------+-------------------------+---------+
|code|count |description              |frequence|
+----+------+-------------------------+---------+
|200 |874321|Succes de la requete     |0.2001   |
|403 |874062|Acces refuse             |0.2001   |
|404 |873765|Ressource non trouvee    |0.2      |
|301 |873596|Redirection              |0.2      |
|500 |873323|Erreur interne du serveur|0.1999   |
+----+------+-------------------------+---------+

Sauvegarde dans la collection 'batch_http_codes'
Traitement termine avec succes.
\end{lstlisting}
    \end{exampleblock}
\end{frame}

\begin{frame}[fragile]{Résultats Batch (codes HTTP)}
    
    \begin{exampleblock}{Document stocké (batch\_http\_codes)}
\begin{lstlisting}[language=Python]
[
  {
    _id: ObjectId('697b70bcf6cb610c2cded206'),
    code: 404,
    count: 873765,
    description: 'Ressource non trouvee',
    frequence: 0.2
  },
  {
    _id: ObjectId('697b70bcf6cb610c2cded207'),
    code: 500,
    count: 873323,
    description: 'Erreur interne du serveur',
    frequence: 0.1999
  }, 
...
\end{lstlisting}
    \end{exampleblock}
\end{frame}

\begin{frame}[fragile]{Résultats Batch (top 10 IPs)}
    
    \begin{exampleblock}{Affichage console (Master Spark)}
\begin{lstlisting}[language=Bash]
Resultats : Top 10 des adresses IP les plus actives
+-------------+-----+
|ip           |count|
+-------------+-----+
|192.168.0.250|4620 |
|192.168.1.52 |4571 |
|192.168.1.243|4558 |
|192.168.0.32 |4553 |
|192.168.2.52 |4534 |
|192.168.0.132|4531 |
|192.168.3.82 |4529 |
|192.168.3.154|4524 |
|192.168.0.101|4522 |
|192.168.1.49 |4521 |
+-------------+-----+

Sauvegarde dans la collection 'batch_top_ips'
\end{lstlisting}
    \end{exampleblock}
\end{frame}

\begin{frame}[fragile]{Résultats Batch (top 10 IPs)}
    
    \begin{exampleblock}{Document stocké (batch\_top\_ips)}
\begin{lstlisting}[language=Python]
[
  {
    _id: ObjectId('697b7107bdbce136d4f6d235'),
    ip: '192.168.0.132',
    count: 4531
  },
  {
    _id: ObjectId('697b7107bdbce136d4f6d237'),
    ip: '192.168.3.82',
    count: 4529
  },
  {
    _id: ObjectId('697b7107bdbce136d4f6d239'),
    ip: '192.168.3.154',
    count: 4524
  },
...
\end{lstlisting}
    \end{exampleblock}
\end{frame}

% --- SECTION 4: STREAMING ---
\section{Traitements Streaming (le présent)}

\begin{frame}{Analyse Streaming : le défi du temps réel}
    \textbf{Génération :} Si un client se connecte au socket TCP, le \texttt{data-generator} lit en boucle le fichier de logs et envoie un a un les logs ayant le même timestamp. Il attend jusuq'au prochain timestamp et recommence jusqu'à déconnexion du client.
    
    \vspace{0.5cm}
    \textbf{Objectif :} Détecter des pics d'erreurs (404/500) \textit{pendant} qu'ils se produisent.
    
    \begin{block}{Problématiques Big Data}
    \begin{enumerate}
        \item \textbf{Event Time :} Utiliser l'heure du log, pas l'heure de la machine.
        \item \textbf{Données en retard :} Gérer les latences réseau.
    \end{enumerate}
    \end{block}
    
    \textbf{Solution :} Utilisation du \textbf{Watermarking} sur les timestamps des logs (tolérance de 1 min).
\end{frame}

\begin{frame}[fragile]{Fenêtrage et agrégation}
    Utilisation de \textbf{Tumbling Windows} (Fenêtres glissantes) de 10s.
    
\begin{lstlisting}[language=Python]
windowed_counts = errors_df \
    .withWatermark("timestamp", "1 minute") \ 
    .groupBy(
        window(col("timestamp"), "10 seconds"),
        col("code")
    ) \
    .count()
\end{lstlisting}

    Exemple : "Entre $15:19:50$ et $15:20:00$ $\rightarrow$ 1372 erreurs 500".
\end{frame}

\begin{frame}[fragile]{Analyse Streaming : méthodologie}

    \begin{enumerate}
        \item \textbf{Connexion} au socket TCP et récupération des données dans un \textbf{DataFrame}.
        \item \textbf{Expression régulière} (Spark SQL) pour extraire le timestamp et le code HTTP.
        \item \textbf{Filtre par code} HTTP (\texttt{"404"}, \texttt{"500"}).
        \item \textbf{Watermarking} d'une minute pour les données en retard.
        \item \textbf{Comptage} en regroupant les timestamps par tranche de 10s (pour tester) et par code.
        \item Output mode : \textbf{UPDATE}.
        \item Sauvegarde dans MongoDB (append).
    \end{enumerate}
 
\end{frame}

\begin{frame}[fragile]{Stratégie d'écriture : UPSERT}
    \textbf{Problème :} Éviter les doublons lors des mises à jour de fenêtres.
    \vspace{0.1cm}
    
    \textbf{Solution :} \textbf{UPSERT} MongoDB avec clé composite. 
    \vspace{0.1cm}
    
    \hspace{1cm}\textit{NB. On reste en UPDATE sur l'output Spark.}
    
\begin{lstlisting}[language=Python]
.mode("append") \
.option("operationType", "replace") \
.option("idFieldList", "start_time,end_time,code")
\end{lstlisting}

    \begin{itemize}
        \item \textbf{Clé unique :} début + fin + code HTTP.
        \item Garantit la consistance des données.
    \end{itemize}
\end{frame}

% --- SECTION 5: CONCLUSION ---
\section{Bilan}

\begin{frame}{Conclusion et perspectives}
    \begin{columns}[t]
        \column{0.5\textwidth}
        \textbf{Bilan Technique :}
        \begin{itemize}
            \item[$\checkmark$] Architecture 100\% conteneurisée.
            \item[$\checkmark$] Pipeline complet : ingestion $\rightarrow$ traitement $\rightarrow$ persistance.
            \item[$\checkmark$] Problèmes de synchronisation et de doublons résolus.
        \end{itemize}

        \column{0.5\textwidth}
        \textbf{Perspectives :}
        \begin{itemize}
            \item Tester sur des \textbf{données et flux réels}.
            \item Remplacer le socket TCP par \textbf{Kafka}.
            \item Connecter \textbf{Grafana} ou \textbf{Metabase} pour faire de la BI.
        \end{itemize}
    \end{columns}
    
    \vspace{1cm}
    \begin{center}
        \Large \textbf{Merci pour votre attention.} \\
    \end{center}
\end{frame}

\end{document}